\section{Análisis del código fuente de LlamaIndex: flujo de Query}

\subsection{Proceso de recuperación}

\begin{enumerate}
    \item Se calcula el Embedding Vector del Query del usuario (también de 1536 dimensiones)
    \item Se calcula la similitud entre el Query Embedding y todos los Node Embedding
    \item Por defecto se toman los \textbf{Top-2} Node más relevantes
\end{enumerate}

El cálculo de similitud admite varias métricas: distancia euclidiana, producto punto, similitud coseno, entre otras.

\subsection{Prompt Template}

\begin{importantbox}{QA Template (extraído del código fuente)}
Plantilla central:
\begin{quote}
\textit{Context information is below. \{context\} Given the context information and not prior knowledge, answer the query. Query: \{query\}}
\end{quote}
Restricción clave: \textbf{responder únicamente con base en la Context Information, sin usar el conocimiento previo del modelo.}
\end{importantbox}

Además existen la Refine Template (que perfecciona una respuesta ya existente a partir de nuevo Context) y la Summarize Template, entre otras.

\subsection{Construcción del mensaje y llamada a la API}

El mensaje final construido incluye:
\begin{enumerate}
    \item \textbf{System Message}: definición de rol — "eres un experto en QA confiable, que siempre responde con base en la Context Information, sin usar conocimiento previo"
    \item \textbf{User Message}: incluye la Context Information (los documentos relevantes Top-K) + el Query
\end{enumerate}

El Response Mode por defecto es Compact (compacto).

\subsection{Resumen de la sección}

Flujo del Query: Query Embedding $\rightarrow$ coincidencia de similitud $\rightarrow$ Top-K Node $\rightarrow$ ensamblaje de la Prompt Template $\rightarrow$ el LLM genera la respuesta. Todo el proceso ocupa solo unas cinco líneas de código del framework, pero el código fuente interno es muy minucioso.

